% Part:sets-functions-relations
% Chapter: size-of-sets
% Section: pairing

\documentclass[../../../include/open-logic-section]{subfiles}

\begin{document}

\olfileid{sfr}{siz}{pai}
\olsection{配对函数与编码}

\begin{explain}
Cantor 的之字形法使得 $\Nat^n$ 的可数性一目了然。
现在让我们聚焦于表示 $\Nat^2$ 的阵列。
沿阵列中的之字形路线计数，我们可以验证 $\tuple{1,2}$ 对应数字~$7$。
然而，如果我们能更直接地计算出这个对应关系，将更为理想。
也就是说，若有之字形枚举的\emph{逆} $g\colon \Nat^2 \to \Nat$，使得
\[
g(\tuple{0,0}) = 0, \;
g(\tuple{0,1}) = 1, \;
g(\tuple{1,0}) = 2, \; \dots,  
g(\tuple{1,2}) = 7, \; \dots
\]
这将使我们能精确计算出 $\tuple{n, m}$ 在该枚举中出现的位置。

事实上，基于两点观察，我们可以直接定义 $g$。第一：若第 $n$ 行、第 $m$ 列的值为~$v$，则第 $(n+1)$ 行、第 $(m-1)$ 列的值为 $v + 1$。第二：该枚举的第一行由三角形数构成，以 $0, 1, 3, 6, \dots$ 开头。第 $k$ 个三角形数是所有小于 $k$ 的自然数之和，可计算为 $k(k+1)/2$。综合这两点观察，考虑如下函数：
\[
  g(n,m) = \frac{(n+m+1)(n+m)}{2} + n
\]
我们通常就写成 $g(n, m)$ 而不是 $g(\tuple{n, m})$，因为这样更为直观。
其含义是先求出第 $(n+m)^\text{th}$ 个三角形数，再加上 $n$。
它恰以我们所希望的方式填充了整个阵列。
特别地，数对 $\tuple{1, 2}$ 被映射到 $\frac{4 \times 3}{2} +
1 = 7$。

函数 $g$ 是某个序对集合之枚举的\emph{逆}。这样的函数称为\emph{配对函数}。
\end{explain}

\begin{defn}[配对函数]
  若函数 $f\colon A \times B \to \Nat$ 是单射的，则称其为算术\emph{配对函数}。我们也说 $f$ \emph{编码}了 $A \times B$，且 $f(x,y)$ 是 $\tuple{x,y}$ 的\emph{编码}。
\end{defn}

\begin{explain}
例如，我们可以用配对函数来编码自然数对；换言之，我们可以用\emph{单个}数来表示每一\emph{对}元素。利用配对函数的逆，我们可以对该数进行\emph{解码}，即找出它所代表的序对。
\end{explain}

\begin{prob}
给出所有非负有理数集合的一个枚举。
\end{prob}

\begin{prob}
证明 $\Rat$ 是可数的。回想一下，任何有理数都可以写成分数 $z/m$ 的形式，其中 $z \in \Int$，$m \in \Nat^+$。
\end{prob}

\begin{prob}
定义 $\Bin^*$ 的一个枚举。
\end{prob}

\begin{prob}
回想初级逻辑课程中的内容：每个可能的真值表都表达一个真值函数。
换言之，真值函数就是从 $\Bin^k$ 到 $\Bin$ 的所有函数（对某个~$k$ 而言）。
证明全体真值函数构成的集合是可数的。
\end{prob}

\begin{prob}
证明任意无限可数集的所有有限子集构成的集合是可数的。
\end{prob}

\begin{prob}
$\Nat$ 的子集称为\emph{共有限的}，当且仅当它是 $\Nat$ 中某个有限集的补集；
也就是说，$A \subseteq \Nat$ 是共有限的当且仅当 $\Nat\setminus A$ 是有限的。
令 $I$ 为由 $\Nat$ 的所有有限子集和共有限子集恰好构成的集合。
证明 $I$ 是可数的。
\end{prob}

\begin{prob}
证明可数个可数集的并仍是可数的。
也就是说，若 $A_1$, $A_2$, \dots{} 均为集合，且每个 $A_i$ 都是可数的，
则它们的并 $\bigcup_{i=1}^\infty
A_i$ 也是可数的。[注意：本题较难！]
\end{prob}

\begin{prob}
设 $f \colon A \times B \to \Nat$ 为任意配对函数。
证明 $f$ 的逆是 $A \times B$ 的枚举。
\end{prob}

\begin{prob}
给出一个能编码 $\Nat^3$ 的具体函数。
\end{prob}

\end{document}